% Header aus der Vorlage
\documentclass[a4paper,11pt, footheight=26pt]{scrartcl}
\usepackage[head=23pt]{geometry}	% head=23pt umgeht Fehlerwarnung, dafür größeres "top" in geometry
\geometry{a4paper, top=30mm, left=20mm, right=20mm, bottom=22mm,headsep=10mm, footskip=12mm}

% Input inkl. Umlaute, Silbentrennung
\usepackage[T1]{fontenc}
\usepackage[utf8]{inputenc}
\usepackage[ngerman]{babel}
\usepackage{csquotes}	% Anführungszeichen

% HTW Corporate Design: Arial (Helvetica)
\usepackage{helvet}
\renewcommand{\familydefault}{\sfdefault}

% Style-Aufhübschung
\usepackage{soul, color}	% Kapitälchen, Unterstrichen, Durchgestrichen usw. im Text
\usepackage{scrlayer-scrpage}	% Kopf-/Fußzeile
%\usepackage{titleref}
\usepackage[perpage]{footmisc}	% Fußnotenzählung Seitenweit, nicht Dokumentenweit

% Mathe usw.
\usepackage{amssymb}
\usepackage[fleqn]{amsmath}	% fleqn: align-Umgebung rechtsbündig
\usepackage{xcolor}
\usepackage{esint}	% Schönere Integrale, \oiint vorhanden
\everymath=\expandafter{\the\everymath\displaystyle}	% Mathe Inhalte werden weniger verkleinert
\usepackage{wasysym}	% mehr Symbole, bspw \lightning

% tikz usw.
\usepackage{tikz}
\usepackage{pgfplots}
\pgfplotsset{compat=1.11}	% Umgeht Fehlermeldung
\usetikzlibrary{graphs}
%\usetikzlibrary{through}	% ???
\usetikzlibrary{arrows}
\usetikzlibrary{arrows.meta}	% Pfeile verändern / vergrößern: \draw[-{>[scale=1.5]}] (-3,5) -> (-3,3);
\usetikzlibrary{automata,positioning} % Zeilenumbruch im Node node[align=center] {Text\\nächste Zeile}
\usetikzlibrary{patterns}	% Schraffierte Füllung
\tikzstyle{reverseclip}=[insert path={	% Inverser Clip \clip
	(current page.north east) --
	(current page.south east) --
	(current page.south west) --
	(current page.north west) --
	(current page.north east)}
% Nutzen: 
%\begin{tikzpicture}[remember picture]
%\begin{scope}
%\begin{pgfinterruptboundingbox}
%\draw [clip] DIE FLÄCHE, IN DER OBJEKT NICHT ERSCHEINEN SOLL [reverseclip];
%\end{pgfinterruptboundingbox}
%\draw DAS OBJEKT;
%\end{scope}
%\end{tikzpicture}
]	% Achtung: dafür muss doppelt kompliert werden!
\usepackage{graphpap}	% Grid für Graphen

% Tabular
\usepackage{longtable}	% Große Tabellen über mehrere Seiten
\usepackage{multirow}	% Multirow/-column: \multirow{2[Anzahl der Zeilen]}{*[Format]}{Test[Inhalt]} oder \multicolumn{7[Anzahl der Reihen]}{|c|[Format]}{Test2[Inhalt]}
\renewcommand{\arraystretch}{1.3} % Tabellenlinien nicht zu dicht
\usepackage{colortbl}
\arrayrulecolor{gray}	% heller Tabellenlinien

% Nützliches
\usepackage{verbatim}	% u.a. zum auskommentieren via \begin{comment} \end{comment}
\usepackage{tabto}	% Tabs: /tab zum nächsten Tab oder /tabto{.5 \CurrentLineWidth} zur Stelle in der Linie
\NumTabs{6}	% Anzahl von Tabs pro Zeile zum springen
\usepackage{listings} % Source-Code mit Tabs
\usepackage{lstautogobble} 
\usepackage{enumitem}	% Anpassung der enumerates
\setlist[enumerate,1]{label=\arabic*.)}	% global andere Enum-Items
\usepackage{letltxmacro} % neue Definiton von Grundbefehlen
% Nutzen:
%\LetLtxMacro{\oldemph}{\emph}
%\renewcommand{\emph}[1]{\oldemph{#1}}

% Einrichtung von lst
\lstset{
basicstyle=\ttfamily, 
mathescape=true, 
escapeinside=||, 
autogobble, 
tabsize=2,
basicstyle=\footnotesize\sffamily\color{black},
frame=single,
rulecolor=\color{lightgray},
numbers=left,
numbersep=5pt,
%basicstyle=\footnotesize\sffamily\color{black},
commentstyle=\color{gray},
%frame=single,
%numbers=left,
%numbersep=5pt,
numberstyle=\tiny\color{gray},
keywordstyle=\color{green},
%showspaces=false,
showstringspaces=false,
stringstyle=\color{orange},
tabsize=2,
literate=%
    {Ö}{{\"O}}1
    {Ä}{{\"A}}1
    {Ü}{{\"U}}1
    {ß}{{\ss}}1
    {ü}{{\"u}}1
    {ä}{{\"a}}1
    {ö}{{\"o}}1
    {~}{{\textasciitilde}}1
}

% BibTeX
\usepackage[backend=bibtex, bibencoding=ascii]{biblatex}	% BibTeX
\usepackage{makeidx}
%\makeglossary
%\makeindex

% Grafiken
\usepackage{graphicx}
\usepackage{epstopdf}	% eps-Vektorgrafiken einfügen

% pdf-Setup
\usepackage{pdfpages}
\usepackage[bookmarks,%
bookmarksopen=false,% Klappt die Bookmarks in Acrobat aus
colorlinks=true,%
linkcolor=black,%
citecolor=red,%
urlcolor=green,%
]{hyperref}

% Titel, Autor usw. werden vor dem Anfang des Dokuments in einem Rutsch definiert…
\newcommand{\DTitel}[1]{\newcommand{\Dokumententitel}{#1}}
\newcommand{\DUntertitel}[1]{\newcommand{\Dokumentenuntertitel}{#1}}
\newcommand{\DAutor}[1]{\newcommand{\Dokumentenautor}{#1}}
\newcommand{\DNotiz}[1]{\newcommand{\Dokumentennotiz}{#1}}
% … Deswegen folgendes erst Nach Dokumentenbeginn ausführen:
\AtBeginDocument{
	\hypersetup{
		pdfauthor={\Dokumentenautor},
		pdftitle={HTW Dresden | \Dokumententitel - \Dokumentenuntertitel},
	}
	\automark[section]{section}
	\automark*[subsection]{subsection}
	\pagestyle{scrheadings}
	\ihead{\includegraphics[height=1.7em]{../LaTeX_master/HTW-Logo.eps}}
	\ohead{\Dokumententitel}
	\ofoot{\footnotesize{\textcolor{darkgray}{Mitschrift von\\ \Dokumentenautor}}}
	% Titelseite
	\title{\includegraphics[width=0.35\textwidth]{../LaTeX_master/HTW-Logo.eps}\\\vspace{0.5em}
	\Huge\textbf{\Dokumententitel} \\\vspace*{0,5cm}
	\Large \Dokumentenuntertitel \\\vspace*{4cm}}
	\author{\textcolor{darkgray}{Mitschrift von \Dokumentenautor} \vspace*{1cm}\\\Dokumentennotiz}
}

%% EIGENE BEFEHLE

%Farbdefinitionen
\definecolor{red}{RGB}{180,0,0}
\definecolor{green}{RGB}{75,160,0}
\definecolor{blue}{RGB}{0,75,200}
\definecolor{orange}{RGB}{255,128,0}
\definecolor{yellow}{RGB}{255,245,0}
\definecolor{purple}{RGB}{75,0,160}
\definecolor{cyan}{RGB}{0,160,160}
\definecolor{brown}{RGB}{120,60,10}

% Textfarbe ändern
\newcommand{\tred}[1]{\textcolor{red}{#1}}
\newcommand{\tgreen}[1]{\textcolor{green}{#1}}
\newcommand{\tblue}[1]{\textcolor{blue}{#1}}
\newcommand{\torange}[1]{\textcolor{orange}{#1}}
\newcommand{\tyellow}[1]{\textcolor{yellow}{#1}}
\newcommand{\tpurple}[1]{\textcolor{purple}{#1}}
\newcommand{\tcyan}[1]{\textcolor{cyan}{#1}}
\newcommand{\tbrown}[1]{\textcolor{brown}{#1}}

% Umstellen der Tabellen Definition
\newcommand{\mpb}[1][.3]{\begin{minipage}{#1\textwidth}\vspace*{3pt}}
\newcommand{\mpe}{\vspace*{3pt}\end{minipage}}

\newcommand{\resultul}[1]{\underline{\underline{#1}}}
\newcommand{\parskp}{$ $\\}	% new line after paragraph
\newcommand{\corr}{\;\widehat{=}\;}
\newcommand{\mdeg}{^{\circ}}

\newcommand{\nok}[2]{\begin{pmatrix}#1\\#2\end{pmatrix}}	% n über k

\bibliography{../Literatur/HTW_Literatur.bib}

% Definition von Titel, Autor usw.
\DTitel{Programmierung II}
\DUntertitel{Vorlesungsskript}
\DAutor{Falk-Jonatan Strube}
\DNotiz{Vorlesung von Dr. Arnold Beck}

\begin{document}

\maketitle
\newpage
\tableofcontents
\newpage

\section*{Hinweise}
Unterlagen unter:
\begin{lstlisting}[language=bash]
cd /home/rex/fi1/nestler/Programmierung_II_2016/
\end{lstlisting}
\paragraph{Compiler} 
\begin{itemize}
\item Intel i16, i13 (für Linux oder Visual Studio) \url{www.hocomputer.de} (kostenpflichtig)
\item gcc 5.3, 4.85 \url{gcc.gnu.org}
\end{itemize}
Zugriff auf Windows-Programme (Visual Studio 2013) in Linux-Laboren:
\begin{lstlisting}[language=bash]
rdesktop -f its56 # oder its59
\end{lstlisting}
Empfohlene Literatur: Breymann\cite{breymann2009c++}

\chapter{C++}
\section{Ein- und Ausgabe}
(siehe Folie CPP\_01\_stdio)
\begin{lstlisting}[language=C++]
#include <iostream>	// alternativ zu <stdio.h> in C
using namespace std;	// namespace: für bestimmte Abkürzungen (bpsw. cout anstatt std::cout)
// Hinweis: "::" zeigt, dass das davorstehende "static" ist (hier: std)

class integer {}; // class: Vergleichbar mit typedef

int main() {
	integer i0;	// i0: Instanz bzw. Objekt der Klasse integer
	cin.get();	// Eingabe abwarten
}
\end{lstlisting}
(vgl. integer.cpp)
\begin{lstlisting}[language=C++]
#include <iostream>
using namespace std;

class integer { // int - Variable in class verpacken
//	private:    // private ist default
		int i;  // this->i bzw. (*this).i
		// private nur für andere Klassen, andere Instanzen der gleichen Klasse können drauf zugreifen
	public: // wenn nichts steht, wird automatisch das vorherige angenommen. Hier: private
		integer(int i=0):i(i){  // Konstruktor und Defaultkonstruktor
			// i=0 default Wert, wenn keiner Angegeben
			// :i(i) übergebenes i wird dem i der Klasse zugewiesen: i_1(i_2) i_1 ist this->i, und i_2 ist das übergebene i
			cout<<"integer-Objekt i = "<<this->i<<endl; 
		}

		int get(){ return i; }

		void set(int i=0){ this->i = i; }

		// statische Methode: aufrufbar ohne Instanziierung der Klasse
		static integer add(integer i1, integer i2){  // Wertkopien von i1 und i2
			// return integer(i1.i + i2.i); // alternativ und explizit: Konstruktor-Aufruf
			// return erstellt eine Kopie (mit Konstruktur erstellt)!
			// return i1.get()+i2.get();    // Umwandlung int nach integer, Aufruf Konstruktor implizit
			return i1.i + i2.i;
			// i1.i Möglich, da innerhalb der Klasse integer und somit privates i sichtbar
		}
};

auto max(int x, int y) -> int { return x>y ? x : y; }	// Lambda-Funktion
// auto: Rückgabetyp ergibt sich aus dem Kontext bzw. über das "-> int"

template<typename Typ1, typename Typ2>	// Weiterentwicklung Makro: wählt automatisch Typ aus
auto quotient(Typ1 a, Typ2 b) -> decltype(a/b) { return a/b; }

auto main() -> int {
	auto k = 0;        // C++11: da 0 vom Typ int ist auch k vom Typ int
	decltype(k) j = 5; // C++11: da k vom Typ int ist auch j vom Typ int
	char *c = nullptr; //C++ 11: Zeigerliteral
	int *ip = NULL;

	integer i0(5), i1=4;     // 2 (alternative) Initialisiierungen von Objekten
	// i1=4 nur möglich, wenn 1 Parameter gefordert ist.
	cout<<"i0.i = "<<i0.get()<<endl;
	// cout im Vergleich zu printf() typsicher.
	cout<<"i0.i + i0.i = "<<integer::add(i0, i0).get()<<endl; // Aufruf static-Methode add
	integer i3 = integer::add(i0, i0);                        // Initialisierung von i3
	cout<<"i3.i     = "<<i3.get()<<endl;
	i0.set(22);
	cout<<"i0.i     = "<<i0.get()<<endl;
	cout<<"max(3,5) = "<<max(3,5)<<endl;
	cout<<"5/3      = "<<quotient(5, 3)<<endl;
	cout<<"5.0/3.0  = "<<quotient(5.0, 3.0)<<endl;
	cout<<"b / 1    = "<<quotient('b', '1')<<endl;
	cin.get();
}
\end{lstlisting}

(vgl iostream.pdf)
\begin{itemize}
\item nach jeder cin Eingabe: „cin.clear();“, damit Fehler ignoriert werden um weiter cin's abhandeln zu können (vgl. robust\_ea)
\end{itemize}
Einlesen:
\begin{lstlisting}[language=C++]
char sc;
cout << "sc=";
cin >> sc;
cin.clear();	// clear, um die Eingabezeile freizumachen, damit man nicht an Falscher Eingabe hängen bleibt
cin.ignore(INT_MAX, '\n');	// braucht #include <limits.h>
cout << "sc" << dec << (int) sc << endl;

// alternativ:
char vb[128];
cout << "s=";
cin.getline(vb, sizeof(vb), '\n');	// lesen als String, dann wieder umwandeln (liest auch Leerzeichen ein)
sc = atoi(vb);	// braucht #include <cstdlib>

// alternative zu getline:
cin.get(...);	// lässt aber \n im Strom
cin.get();

// alternativ
cin >> setw(sizeof(vb)) >> vb;	// verhindert Überlauf
sc = atoi(vb);

// alternativ
String s;	// braucht #include <string>
size_t ie=0;
cin >> s;
unsigned int ni = stoi(s, &ie, 10);

// alternativ
getline(cin, s, '\n');
double d = stod(s, &ie);

// zum compilieren: g++ p2a1.cpp -std=c++11 -o a.out
\end{lstlisting}

robust\_ea1.cpp:
\begin{lstlisting}[language=C++]
	do {
		cout<<"d = "; cin>>d; 	// einlesen
		if(cin.eof()) break;		// break bei Strg+D oä.
		if(cin.fail() || (cin.peek() != '\n')){		// ist nächstes Zeichen ungültig?
			cin.clear(); cin.ignore(INT_MAX,'\n');	// Strom zurücksetzen und zum \n gehen
			continue;
		}
		break; % Schleife verlassen, wenn korrekte Eingabe
	} while(true);
	
	if(cin.eof()){ cin.clear(); cout<<"eof\n"; }
	else {
		cin.clear(); cin.ignore(INT_MAX,'\n');
		cout<<"Wert d = "<<d<<endl;
	}
	cin.ignore();
\end{lstlisting}

\section{Defaultargumente}
Defaultargumente müssen immer von rechts angefangen definiert sein:
\begin{lstlisting}[language=C++]
myFunc(int i = 5, int j = 7)	// korrekt
myFunc(int i, int j = 7)			// korrekt
myFunc(int i = 5, int j)			// falsch!!!
\end{lstlisting}

\section{Überladen}
overload.pdf

Hinweis: cast auf zwei Möglichkeiten:
\begin{lstlisting}[language=C++]
int i = 5;
double d;
d = (long) i;
d = long(i);
\end{lstlisting}

\section{Typisierte Konstanten}

\section{Referenzen}
referenzen.pdf\\
Ein Speicherplatz wird mit mehreren Variablen-Namen beschrieben.

\section{String0}
\section{String}
\section{const}
const\_mutable.pdf\\
Faustregel: Alle Funktionen, die nichts ändern, immer das const anfügen.
\section{Folge}
Initfolge/Initfolge.pdf
\section{Kopierkonstrukturen}
Kopierkonstruktoren.pdf

\section{Friend}
Zugriff auf private Klassen von außerhalb:
\begin{itemize}
\item Funktionen
\item Alle Methoden anderer Klassen
\item spezifische Methoden anderer Klassen
\end{itemize}
Muss von Klasse mit privater Funktion festgelegt werden. Freundschaft wird nicht „erwidert“, ist nicht reflexiv (bsp. friendreflex).

\section{Vererbung}
vererben.pdf

\section{Virtuelle Funktionen}
\begin{itemize}
\item überschriebene Funktionen 
\item virtuelle Funktionen $\Rightarrow$Späte Bindung
\item rein virtuelle Funktionen\\
eine Klasse mit rein virtuellen Funktionen ist abstrakt
\item VMT virtual method table ABB 124
\end{itemize}
besch1.cpp (Beschäftigte)
\begin{lstlisting}[language=C++]
#include <iostream>
#include <string>
using namespace std;

class Besch{
	private:
		string name;
	public:	
		Besch(string name): name(name){}	// :name(name) --> Member-Initialisierer
		string getName() {return name;}
		void setName(string name) {this->name = name;}
		virtual void display(ostream& os) {os<<name;}	// virtuelle Funktion
		virtual double calc(){return 0.0;}	// normalerweise Geld nicht mit double, sondern ganzzahlig rechnen!
		// wenn Funktion eigentlich nicht benötigt wird: Polymorphes Interface:
		// virtual double calc() = 0; // rein virtuelle Funktion
		// diese rein virtuelle Funktion macht den Datentyp Besch abstract! D.h. davon kann keine Instanz erzeugt werden (wie in der Main passiert)!
};
ostream& operator<<(ostream& os, Besch& b) {b.display(os);}

class Arbeiter:public Besch{	// public: durchsichtige Vererbung
	private:
		int stunden;
		double stdLohn;
	public:
		Arbeiter(string name, double stdLohn):Besch(name){ // :Besch(name) --> Basis-Initialisierer
			this->stdLohn = stdLohn;
		}
		int getStunden() { return stunden; }
		void setStunden(int stunden) { this->stunden = stunden; }	// usw. getter und setter
		void display(ostream& os){ Besch::display(); os<<""<<stdLohn<<" "<<stunden<<" "<<calc();}
		double calc(){ return stdLohn * stunden; }
};
ostream& operator<<(ostream& os, Arbeiter& a) { a.display(os); }

class Angest:public Besch{
	private:
		double gehalt;
	public:
		Angest(string name, double stdLohn):Besch(name){
			this->gehalt = gehalt;
		}
		void display(ostream& os){ Besch::display()<<" "; os<<gehalt<<" ";}
		double calc(){ return gehalt; }
};
ostream& operator<<(ostream& os, Angest& a) { a.display(os); }

class Haendler:public Arbeiter{
	private:
		double prov;
		double ums;
	public:
		Haendler(string name, double stdLohn, double prov):Arbeiter(name, stdLohn){this->prov=prov;};
		void setUms(double ums){this->ums = ums;}
		double calc(){return Arbeiter::calc() ++ ums*prov;}
		void display(ostream &os){Arbeiter::display(os);}
}

int main() {
	Besch* b1 = new Besch("Hans Huckebein");
	b1->display(cout); cout<< " Euro"<<endl;
	cout<< (*b1) << " Euro" << endl;
	cout<< "=============="<<endl;
	Arbeiter* b2 = new Arbeiter("Moriz Lehmann",  8.99);
	b2->setStunden(140);
	b2->display(cout); cout<< " Euro"<<endl;
	cout<< (*b2) << " Euro"<< endl;
	cout<< "=============="<<endl;
	Angest* b3 = new Arbeiter("Friedrich Lempel",  3099.00);
	b3->display(cout); cout<< "Euro"<<endl;
	cout<< (*b3) << " Euro"<< endl;
	cout<< "=============="<<endl
	Haendler* b4 = new Handler("Bang Ohlufson",  15.80, 0.1);
	b4->setStunden(350);
	b4->setUms(12000);
	b4->display(cout); cout<< "Euro"<<endl;
	cout<< (*b3) << " Euro"<< endl;
	cout<< "=============="<<endl
	
	Besch* be[] = {b1,b2,b3,b4};
	for (int i=0; i<4; i++){
		cout<<i<<": "; 
		be[i]->display(cout); cout<<endl;	// gibt 0 für alle zurück, da nur calc() (und display()) vom Typ Besch ausgeführt, da das Array aus Besch besteht! Die Funktion calc() wurde überschrieben (Vgl. Überladen, wenn bei gleichen Namen verschiedene Paramater angegeben sind)!
		// Pointer im Array bestimmt die ausgeführten Funktionen!
		// Lösung: "virtual" als Schlüsselwort vor Funktion (s.o.). Angeben bei erstem Auftauchen in Vererbungshierarchie
		cout<<(*(be[i]))<<endl<<"+++++++++++++++"<<endl;	// geht auch, da im Operator (die ja nicht virtuell gemacht werden kann, weil sie kein Member ist) die virtuelle display-Fkt aufgerufen wird.
	}	
	
	return 0;
}
\end{lstlisting}
Vererbungs Hierarchie:\\
ABB 125\\
Potentielle Mehrfachvererbung (bei Manager):
\begin{itemize}
\item Hat mehrfache Namen (wegen unterschiedlichen Vererbungslinien)
\end{itemize}
Problemlösung: Arbeiter, Angest und Händler von virtual Besch erben lassen (dort, wo es sich aufzweigt):
\begin{lstlisting}[language=C++]
class Arbeiter: virtual public Besch{};
\end{lstlisting}
Hat immer noch Einschränkungen!

\section{Templates}
Templates Folie (enthält auch exception).
Vor allem für Containerklassen/Collections verwendet (Vektoren, Hashtables, Listen).\\
Schafft Möglichkeit Klassen parameterbehaftet zu generieren (generische Klassen in Java).\\
Funktionen von Templates müssen (im Gegensatz zur bisherigen Vorgehensweise) im header-File ausprogrammiert werden. Alles, was zum Template gehört kommt ins header-file!!!
\paragraph{Beispiel Array} Intarr.h
\begin{lstlisting}[language=C++]
const int SizeArr = 24;

// Konstruktor, Kopierkonstruktor und Destruktor werden bei Klassen mit Pointer immer benötigt!
class IntArr{
	public:
		IntArr (int Gr = SizeArr ) 	// Konstruktor
		IntArr (const IntArr&) 			// Kopierkonstruktor
		~IntArr() (delete IA;}			// Destruktor
		IntArr & operator = (const IntArr&);
		int& operator [] (int i);	
		int getNum() const {return Size;}
		void resize(int NewSz);
	private:
		int Size;
		int * IA;
};
\end{lstlisting}
Auszug Intarr.cpp
\begin{lstlisting}[language=C++]
...
IntArr :: IntArr (int Sz){
	Size = Sz;
	IA = new int [Size];
	for (int i=0; i<Sz; ++i) IA[i] = 0;
}
\end{lstlisting}
Modifikation um es generisch zu machen (Template), Inhalte von Intarr.cpp müssen eingefügt werden:\\
Intarr.h
\begin{lstlisting}[language=C++]
const int SizeArr = 24;

template <class T>
class IntArr{
	public:
		IntArr (int Gr = SizeArr );
		IntArr (const IntArr<T>&);
		~IntArr() (delete IA;}
		IntArr & operator = (const IntArr<T>&);
		T& operator [] (int i);	
		int getNum() const {return Size;}
		void resize(int NewSz);
	private:
		int Size;
		T * IA;
};

//... + Inhalt von Intarr.cpp!!!
template <class T>
IntArr<T> :: IntArr (int Sz){
	Size = Sz;
	IA = new T [Size];
	for (int i=0; i<Sz; ++i) IA[i] = 0;
}

template <class T>
IntArr<T> :: IntArr (const IntArr<T>& Other){
	Size = Other.Size;
	IA = new T[Size];
	for (int i=0; i<Size; i++) IA[i]=Other.IA[i];
}
usw.
\end{lstlisting}
in der Main dann:
\begin{lstlisting}[language=C++]
// vorher: IntArr IA (10);
IntArr<int> IA (10);
// oder auch 
IntArr<double> IA (10);
// auch möglich 
IntArr<Fraction> IA (10); // wobei Fraction in einer anderen Klasse definiert ist.
\end{lstlisting}

\paragraph{Beispiel Listen}
cobject.h

\chapter{Java}
\section{Grundlagen}
Java0intro.pdf\\
Klassenname und Dateiname müssen exakt übereinstimmen. Jede Datei darf nur eine Klasse enthalten.
\begin{lstlisting}[language=Java]
public class HelloOne{	// Dateiname: HelloOne.java
	public static void main(String args[]}{	// im Vergleich zu C++ gibt es keinen public-Bereich, die Eigenschaft muss vor jede Funktion/Variabel geschrieben werden.
		System.out.println("welcome to java");
		System.out.println("today is " + new java.util.Date());
		System.out.printf("JAVA test: %d\n", 26);
	}
}
\end{lstlisting}
javac HelloOne.java zum Übersetzen\\
java HelloOne zum Ausführen
\begin{lstlisting}[language=Java]
public class HelloEcho{
	public static void main(String args{]){
		for (int i = 0; i<args.length; i++)
			System.out.println(args[i]);
		// oder:
		for (String s:args)
			System.out.println(s);
		System.exit(0);
	}
}
\end{lstlisting}
\subsection{Parsen von Daten}
\begin{lstlisting}[language=Java]
public class HelloParse{
	public static void main(String args{]){
		int i;
		double d;
		i = Integer.parseInt(args[0]);
		d = Double.parseDouble(args[1});
		System.out.println("i: "+i+" d: "+d);
	}
}
\end{lstlisting}
\subsection{GUI mit AWT}
\begin{lstlisting}[language=Java]
import java.awt.*;
import java.awt.event.*;
import java.util.*;
public class HelloTwo extends Panel {	// extends: erweitert. Also: "erbt von"
	Label l1=new label("welcome to java", Label.CENTER);
	Label l2=new Label("today is "+new Date(), Label.CENTER);	// wegen import javo.util.* kann hier nur Date() geschrieben werden.
	Button close = new Button ("Close");
	
	public HellowTwo(){
		setFont(new Font("System", Font.PLAIN, 24));
		setLayout(new GridLayout(3,1,0,10));
		add(l1);	// alternativ: this.add(l1);
		add(l2);
		close.addActionListener(new ActionListener(){
			public void actionPerformed(ActionEvent e){
				// bpsw erweiterbar mit System.out.printl("Close...");
				System.exit(0);
			}
		});
	}
	
	public static void main(String args[]){
		Frame f=new Frame("first Desktop Application");
		HelloTwo h = new HelloTwo();
		f.add(h);
		f.setVisible(true);
		f.pack();
	}
}
\end{lstlisting}

\subsection{Applet}
\begin{lstlisting}[language=Java]
import java.applet.*;
import java.awt.*;

public class HelloThree extends Applet{
	public void init(){
		Panel h= new HelloTwo();
		add(h);
	}
}
\end{lstlisting}
\begin{lstlisting}[language=HTML]
<html>
<body>
<applet code=HelloThree width=940 height=540></applet>
</body>
</html>
\end{lstlisting}
Zu jedem Applet gehört ein dazugehöriges HTML-File. Damit kann das Applet über einen Applet-Viewer benutzt werden. Hier funktioniert der Close-Button aber nicht. 
\paragraph{Angepasstes HelloTwo} Hiermit wird der Close-Button implementiert, sodass er im Applet nicht erscheint.
\begin{lstlisting}[language=Java]
// Quit-Button wird erst in main-Funktion eingeführt
\end{lstlisting}
java1.pdf
\subsection{Operatoren}
\begin{itemize}
\item $>>$ schiebt nach rechts und fügt abhängig vom ersten Bit $0$en (1. Bit $0$) oder $1$en (1. Bit $1$) nach. $>>>$ schiebt immer $0$en nach.
\end{itemize}


\newpage
\printbibliography
\end{document}